Contextual Retrievalやmoderation APIのような更新は、model capabilityとは別に、回答を外部情報へgroundingしたり、入出力を分類・制御したりするlayerが製品設計上の主要要素になったことを示す。retrieval改善値やmoderation性能は評価設定に依存するため、release factと性能claimを分離して扱う。\autocite{src-308d34e54e69,src-33d811167435}


factualityやagent/tool executionを測る研究、prompt injectionや外部入力を介した脅威を扱う研究は、能力が高いだけでは安全なexecution stackにならないことを示した。各研究は独自のthreat modelとevaluation setupを持つため、個別結果を『実運用の一般的安全性』へ拡張しない。\autocite{src-e7b786f4a622,src-f96e45232e1f}


alignment fakingやdeliberative alignmentなどの安全性研究は、モデルがより長く考え、より多くの外部作用を持つ時代に、behavior evaluationとcontrol methodを独立した研究対象として扱う必要を示す。ただしcontrolled experimentで観測されたbehaviorを、deployed model全体の恒常的挙動として一般化しない。\autocite{src-3278a129a17e,src-9cf289b414fd,src-e92e95daf8e0}


